\subsection{波兰空间}

定义 1. 拓扑空间 X 称为波兰空间，如果它可可数型度量化（§ 2，第 8 号），且存在一个与 X 的拓扑相容的度量，X 关于它是完备的。

命题 1. a) 波兰空间的每个闭子空间都是波兰空间。

\begin{enumerate}
\def\labelenumi{\alph{enumi})}
\setcounter{enumi}{1}
\item
  可数族波兰空间的积是波兰空间。
\item
  可数族波兰空间的和是波兰空间。
\end{enumerate}

可数型可度量化空间的每个子空间都可数型度量化，完备空间的每个闭子空间都完备（第 II 章，§ 3，第 4 号，命题 8）。可数型可度量化空间的每个可数积仍可数型度量化（§ 2，第 8 号），且完备度量空间的每个可数积关于一个与其拓扑相容的度量是完备度量空间（第 II 章，§ 3，第 5 号，命题 10 及第 IX 章，§ 2，第 4 号，定理 1 的推论 2）。最后，设 \((\mathbf{X}_n)\) 是非空波兰空间的序列，考察积空间 \(\mathbf{Y} = \mathbf{N} \times \prod_{n} \mathbf{X}_n\) ，其中 \(\mathbf{N}\) 赋有离散拓扑；由已证的结果，\(\mathbf{Y}\) 是波兰空间。另一方面，设 \(a_n\) 是 \(\mathbf{X}_n\) 的一点（对每个 \(n\) ），并设 \(f_n\) 是 \(\mathbf{X}_n\) 到 \(\mathbf{Y}\) 中的映射，使得对每个 \(x \in \mathbf{X}_n\) 有

\[
f _ {n} (x) = (n, \left(y _ {p}\right)),
\]

其中 \(y_{p} = a_{p}\) （若 \(p \neq n\) ），而 \(y_{n} = x\) 。若 \(X\) 是诸 \(X_{n}\) 的拓扑和，则显然这样的映射 \(f\) —— \(X\) 到 \(Y\) 中的、与 \(f_{n}\) 在每个 \(X_{n}\) 上（对每个 \(n\) ）一致的映射 —— 是 \(X\) 到 \(f(X)\) 上的同胚；此外，对每个 \(n\) ，\(f_{n}(X_{n})\) 在 \(Y\) 中闭，且族 \((f_{n}(X_{n}))\) 局部有限，因为 \(N\) 是离散的；因此 \(f(X) = \bigcup_{n} f_{n}(X_{n})\) 在 \(Y\) 中闭（第 I 章，§ 1，第 5 号，命题 4），于是由 a) ，\(f(X)\) 是波兰空间。

命题 2. 波兰空间的每个开子空间都是波兰空间。

设 X 是波兰空间，d 是 X 上与其拓扑相容的度量，\(U \neq X\) 是 X 的开子集。设 V 是 \(R \times X\) 中由满足下列条件的点 \((t, x)\) 组成的子集：

\[
t. d (x, \mathrm{X} - \mathrm{U}) = \mathrm{I};
\]

由 § 2 第 2 号命题 3 ，V 是闭的，因而是波兰空间（命题 1）。由于投影 \(\mathbf{pr}_{2}:\mathbf{R}\times\mathbf{X}\to\mathbf{X}\) 在 V 上的限制是 V 到 U 上的同胚（§ 2，第 2 号，命题 3），故 U 是 X 的波兰子空间。

推论. 每个局部紧、σ 紧、可度量化的空间 X 都是波兰空间。

设 \(X'\) 是由 \(X\) 添加一个无穷远点而得的紧空间；\(X'\) 可度量化且可数型（§ 2，第 9 号，命题 16 的推论），且 \(X'\) 关于其唯一的一致结构是完备的（第 II 章，§ 4，第 1 号，定理 1）。故 \(X'\) 是波兰空间；由于 \(X\) 在 \(X'\) 中开，由此得出 \(X\) 是波兰空间。

命题 3. 设 X 是豪斯多夫拓扑空间。则 X 的波兰子空间的序列 \((\mathbf{A}_{n})\) 的交是波兰子空间。

设 \(f\) 是 \(\mathbf{X}\) 到 \(\mathbf{X}^{\mathbf{N}}\) 中的对角映射{[}集合论，第 II 章，§ 5，第 3 号；回忆 \(f(x) = (y_n)\) ，其中 \(y_n = x\) 对一切 \(n \in \mathbb{N}\) 成立 {]}。我们将用到下述引理：

引理 1. 设 \((\mathbf{A}_{n})\) 是豪斯多夫拓扑空间 X 的子集的序列。则对角映射 \(f: X \to X^{N}\) 在 X 的子空间 \(\bigcap_{n} A_{n}\) 上的限制是 \(\bigcap_{n} A_{n}\) 到 \(\prod_{n} A_{n}\) 的一个闭子空间上的同胚。

此像是 \(\prod_{n} A_n\) 与对角集 \(\Delta = f(X)\) 的交，它在 \(X^N\) 中闭，因为 \(X\) 是豪斯多夫空间（第 I 章，§ 8，第 1 号）；且 \(f\) 是 \(X\) 到 \(\Delta\) 上的同胚。

在命题 3 的假设下，\(\prod_{n} A_{n}\) 是波兰空间（命题 1），故由引理 1 与命题 1 ，\(\bigcap_{n} A_{n}\) 是波兰空间。

推论. 无理数集赋以由实直线 R 诱导的拓扑后是波兰空间。

它是 R 中一个可数开集族的交，这些开集是由单个有理数组成的集的余集。

定理 1. 波兰空间 X 的子空间 Y 是波兰空间，当且仅当 Y 是 X 中可数开集族的交。

该条件的充分性由命题 2 与命题 3 立即得出。为证必要性，设 \(d\) 是与 Y 的拓扑相容且 Y 关于它完备的度量。设 \(\overline{\mathbf{Y}}\) 是 Y 在 X 中的闭包。对每个整数 \(n>0\) ，令 \(\mathbf{Y}_{n}\) 为 \(x\in\overline{\mathbf{Y}}\) 中具有开邻域 U 使 \(\mathbf{U}\cap\mathbf{Y}\) （关于度量 \(d\) ）的直径 \(\leqslant1/n\) 的一切点所成的集。显然 \(\mathbf{Y}_{n}\) 在 \(\overline{\mathbf{Y}}\) 中开且包含 Y 。设 \(x\in\bigcap_{n}\mathbf{Y}_{n}\) ；则 \(x\in\overline{\mathbf{Y}}\) ，且 \(x\) 在 X 中的邻域滤子在 Y 上的迹是柯西滤子（关于 \(d\) ）；故该滤子收敛于 Y 的一点，从而 \(x\in\mathbf{Y}\) 。于是 \(\mathbf{Y}=\bigcap_{n}\mathbf{Y}_{n}\) 。对每个 \(n\) ，设 \(\mathbf{H}_{n}\) 是 X 的开子集，使 \(\mathbf{H}_{n}\cap\overline{\mathbf{Y}}=\mathbf{Y}_{n}\) ，并设 \((\mathbf{U}_{m})\) 是 X 的开子集的序列，使 \(\overline{\mathbf{Y}}=\bigcap_{m}\mathbf{U}_{m}\) （ \(\S\) 2，第 5 号，命题 7）；于是 Y 是可数开集族 \((\mathbf{H}_{n}\cap\mathbf{U}_{m})\) 的交。

推论 1. 空间 X 是波兰空间，当且仅当它同胚于方体 IN 中开集的可数交，其中 I 是 R 的区间{[}O, I{]}。

该条件显然充分，且它是必要的，因为每个可数型可度量化空间都同胚于 \(I^{N}\) 的一个子空间（ \(\S\) 2，第 8 号，命题 12）。

推论 2. 设 X 与 Y 是两个波兰空间，\(f: \mathbf{X} \to \mathbf{Y}\) 是连续映射。若 Z 是 Y 的波兰子空间，则 \(\overline{f}^{1}(Z)\) 是 X 的波兰子空间。

因为 \(Z = \bigcap_{n} Z_n\) ，其中诸 \(Z_n\) 是 \(Y\) 的开子集；故

\[
\overline {{{f}}} ^ {1} (Z) = \bigcap_ {n} \overline {{{f}}} ^ {1} (Z _ {n}),
\]

且诸集 \(\overline{f}^1 (Z_n)\) 在 \(\mathbf{X}\) 中开。

